\begin{textsample}{POS Dim 1 – pr – Score 97.00 – pr/pr\_em\_50.txt}  \label{ex:f1_pos_024}
4.3.2 \textbf{Encaminhamentos} \textbf{metodológicos} Integrando os demais componentes da Área de Linguagens e suas Tecnologias, esse referencial curricular de Língua \textbf{Inglesa} ( LI ), por meio dos \textbf{encaminhamentos} pedagógicos, busca \textbf{subsidiar} o professor em suas práticas de linguagem ( leitura, oralidade e escrita ) ; “ práticas de \textbf{compreensão} e produção escrita ; práticas de \textbf{compreensão} e produção oral ; conscientização linguística e conscientização \textbf{intercultural} ”, sob a perspectiva trans e interdisciplinar e contextualizada, de modo que possa “ garantir aos estudantes do Ensino Médio o desenvolvimento de \textbf{competências} específicas e habilidades a serem alcançadas nessa etapa ” ( Brasil, 2018, p. 481 ).
Destacam -se nesses \textbf{encaminhamentos} os pressupostos fundantes da concepção de ensino da Língua \textbf{Inglesa} como língua franca, e as premissas que \textbf{permeiam} as práticas de linguagem.
A primeira premissa \textbf{diz} respeito à consideração da Língua \textbf{Inglesa} — \textbf{Inglês} como Língua Franca ( ILF ) — no presente contexto mundial, o que \textbf{implica} na \textbf{abordagem} de ensino da LI desvinculada da relação de que serve somente aos falantes nativos de países anglófonos, pois ela deve “ servir como meio de comunicação para indivíduos de diferentes nacionalidades com diferentes repertórios linguístico-culturais ” ( Coutinho, 2017, p. 5 ).
Essa consideração traz em si a marca de mudanças na concepção de correção e níveis de \textbf{proficiência} a serem aceitos no ensino, deslocando conceitos de modelos \textbf{ideais} de falantes da LI ( \textbf{Inglês} Britânico, Americano, Australiano ) e de culturas específicas a serem tratadas na relação língua e cultura na composição dos currículos de LI, definindo, assim, a segunda premissa, que favorece uma educação linguística, em que há a sensibilidade e \textbf{compreensão} de como as línguas e culturas são manifestadas na produção de conhecimentos e visões de mundo.
A terceira premissa destaca os \textbf{letramentos} e \textbf{multiletramentos}, envolvendo as múltiplas linguagens ( verbal, visual, corporal, audiovisual ), que \textbf{impactam} nos processos de significação e nos aspectos cognitivos, afetivos e sociais, integrando diferentes semioses, em processos contínuos de significação, dialogia e ideologias presentes nos textos, com destaque à cultura \textbf{digital} nas práticas sociais, especialmente na cultura juvenil.
A quarta premissa recai sobre “ a importância da \textbf{interdisciplinaridade} ” e o papel da língua \textbf{inglesa} no acesso às diferentes áreas de conhecimento, em diferentes ambientes ( \textbf{virtuais} e não \textbf{virtuais} ) e esferas de circulação social. ( Coutinho, 2017 ).
Diante dos pressupostos e conceitos-chave que \textbf{permeiam} o ensino da LI, no presente contexto, considera -se desconstruir a noção de ensino com a \textbf{abordagem} estruturada das práticas de linguagem ( leitura, oralidade e escrita ).
Embora esse documento apresente os \textbf{quadros} \textbf{organizadores} para as práticas de linguagens, estes se constituem na concepção de organização de currículo, a visão macro de um arcabouço de possibilidades de articulação e integração dos objetos de conhecimento e as sugestões de conteúdos, com \textbf{vistas} ao desenvolvimento de habilidades para a construção de \textbf{competências} específicas na área de linguagens, alinhadas ao desenvolvimento das \textbf{Competências} Gerais ( CG ).
Para \textbf{tanto}, as práticas de linguagem se realizam de \textbf{maneira} integrada na \textbf{abordagem} de \textbf{ensino-aprendizagem} da LI, como prática social, em sua função de linguagem : língua de interação entre sujeitos, em situações reais e situadas de interlocução.
Estudos e pareceres se manifestam em favor dessa ressalva para o ensino da LI no Ensino Fundamental, onde a LI está organizada em eixos, destacando que “ os eixos estão intrinsecamente ligados nas práticas sociais em que estão inseridos ” ( Coutinho, 2017, p.5 ), e que embora a orientação da BNCC se apresente de forma estruturada nos \textbf{quadros}, estes não devem ser compreendidos sob a concepção antiga de ensino fragmentado, e com o \textbf{foco} nas estruturas linguísticas e gramaticais da língua.
Assim considerando, acatamos a coerência da consideração das práticas de linguagem como “ práticas de \textbf{compreensão} e produção oral ”, “ práticas de \textbf{compreensão} e produção escrita ”, “ conscientização linguística ” e “ conscientização \textbf{intercultural} ”, respeitando as características de flexibilidade e integralidade na construção de conhecimentos durante os percursos educativos, “ o estímulo a sua aplicação na vida real, a importância do contexto para dar sentido ao que se aprende e o protagonismo do estudante em sua aprendizagem (... ) ” no Novo Ensino Médio ( Brasil, 2018, p.15 ).
Acrescente -se, aqui, o papel da Proposta Político-Pedagógica da instituição de ensino que manifestará suas escolhas em relação às metodologias e \textbf{abordagens} para o ensino da língua, seleção de gêneros orais ou escritos, projetos e espaços de aprendizagem da LI em seu contexto local, \textbf{reiterando} as decisões e orientações sobre a Base Nacional Comum Curricular e currículos : > [... ] a BNCC e currículos têm papéis complementares para assegurar as aprendizagens essenciais definidas para cada etapa da Educação Básica, [... ].
São essas decisões que vão adequar as proposições da BNCC à realidade local, considerando a autonomia dos sistemas ou das redes de ensino e das instituições escolares, como também o contexto e as características dos \textbf{alunos}.
Essas decisões, [... ] referem -se, entre outras ações, a ) contextualizar os conteúdos dos componentes curriculares, identificando estratégias para apresentá -los, representá -los, exemplificá -los, conectá -los e torná -los significativos, com base na realidade do lugar e do tempo nos quais as aprendizagens estão situadas, ( Brasil, 2018, p. 16-17, \textbf{grifos} próprios ).
O presente referencial curricular, nesse contexto, tem o objetivo de apresentar pressupostos e procedimentos que possam se constituir em contribuições para práticas de linguagem significativas e alinhadas aos pressupostos fundantes orientados pela BNCC ( 2018 ).
As orientações para as práticas \textbf{discursivas} no componente LI, que serão apresentadas a seguir, trazem a \textbf{descrição} das possibilidades curriculares e procedimentos didático-pedagógicos, com \textbf{vistas} ao desenvolvimento das habilidades da BNCC ( 2018 ), vinculadas às \textbf{competências} específicas da Área de Linguagens e suas Tecnologias, integrando -se por sua vez aos componentes da Área, em relações trans e interdisciplinares contextualizadas e em consonância com as \textbf{competências} gerais para a Formação Geral Básica ( FGB ).
As práticas \textbf{discursivas} em leitura, oralidade e escrita em LI integram -se na Área para o desenvolvimento da habilidade EM13LGG101, pelos procedimentos comuns de \textbf{análise} e exploração de seus objetos de estudo.
Dessa forma, a leitura, escuta, apreciação, experimentação, \textbf{análise} de discursos e atos de linguagem podem ocorrer por meio de \textbf{contextualização}, \textbf{problematização} e propostas de investigação que permitam ao estudante relacionar aspectos contextuais e escolhas de recursos expressivos ( linguísticos, gestuais, artísticos, \textbf{multissemióticos} ), na produção de sentidos e na revelação de intencionalidades, de modo que os jovens possam fazer escolhas fundamentadas de discursos e de atos de linguagem, de acordo com interesses pessoais e coletivos, seus valores e seus projetos de vida.
Ao trabalhar esta habilidade, a área contribui para o desenvolvimento do pensamento autônomo, da responsabilidade e do uso das linguagens assentado em princípios éticos, democráticos, inclusivos, sustentáveis e solidários, conforme previsto na \textbf{Competência} Geral 10.
A habilidade EM13LGG102 integra -se pela exploração conjunta de temáticas \textbf{complexas} ( racismo, estereótipos de gênero, feminismo, violência, privacidade nas mídias sociais ), evidenciando -se na ampliação do repertório \textbf{crítico} dos jovens e na diversidade de modos de exprimir sua visão \textbf{crítica} sobre esses ou outros temas que lhes sejam significativos.
É possível, ainda, a integração com a área Ciências Humanas e Sociais Aplicadas por meio de procedimentos comuns, pois a \textbf{análise} de discursos e atos de linguagem, com procedimentos analíticos, permite que o estudante construa ferramentas necessárias para reconhecer a \textbf{pluralidade} de visões de mundo e interesses que se \textbf{materializam} em textos e para lidar com a complexidade de discursos sobre temas, questões e conflitos \textbf{contemporâneos}.
Dessa forma, os atos e as práticas de linguagem trabalhados na área colocam -se a serviço da ampliação cultural do estudante, do desenvolvimento da \textbf{capacidade} argumentativa e do compromisso com a cooperação e o respeito às diferenças, como previsto nas \textbf{Competências} Globais 1, 3 e 9.
A habilidade EM13LGG103 integra -se dentro da própria área pela adoção de metodologias ativas para estudo das semioses, favorece a agência dos jovens, que passam a se \textbf{implicar} no processo de aprendizagem.
Ao estudar uma determinada semiose ( as materialidades e linguagens em processos da criação teatral ; marcas de tempo nas desinências verbais e suas possibilidades expressivas em textos de gêneros com predomínio do narrar ; conjunções na \textbf{coesão} e na coerência em textos de gêneros com predomínio do \textbf{argumentar} ; a constituição do gesto em diferentes práticas corporais e contextos, etc. ), o estudante poderá analisar usos em diferentes contextos e práticas e sintetizar suas descobertas de diferentes formas.
Em colaboração com outros colegas, poderá : aprofundar e \textbf{aprimorar} seus registros ; resgatar e \textbf{sistematizar} conhecimentos que vêm sendo construídos desde a etapa de escolarização anterior, assim como ampliá -los, com reflexão sobre a \textbf{progressão} nessas aprendizagens.
São metodologias favoráveis para essas práticas : sala de \textbf{aula} invertida, trabalho em ilhas e rotações.
Esse trabalho da área poderá favorecer também o desenvolvimento da \textbf{Competência} Geral 4, que prevê a ampliação do uso das linguagens, para promover o entendimento mútuo em busca da construção de uma sociedade mais justa e equitativa.
Para favorecer o desenvolvimento da habilidade EM13LGG104 de forma significativa, a área poderá propor momentos de culminância, resultado de \textbf{sequências} de atividades ou de projetos, com produção de textos e atos de linguagem diversos, em \textbf{eventos} de \textbf{multiletramentos}, arte e práticas da cultura corporal ( saraus, cinedebate, campeonatos, mostras de artes ), que possam acionar a participação dos jovens com protagonismo, \textbf{mobilizando} saberes e fazeres dos componentes, no planejamento, execução e avaliação de vivências de linguagens voltadas para a comunidade escolar, promovendo a construção de um patrimônio \textbf{artístico-cultural} local.
Para favorecer o desenvolvimento da habilidade EM13LGG105, o componente LI pode promover propostas que deem aos jovens protagonismo na identificação de situações em que queiram intervir, sejam locais ou globais, por meio de projetos com gêneros próprios das culturas juvenis e da cultura de convergência : as reportagens multimidiáticas ( que podem ser criadas pela produção conjunta de uma reportagem, um documentário, \textbf{entrevistas} em áudio ), campanhas de conscientização multimidiáticas e flashmobs integrados.
O trabalho com esta habilidade favorece que o estudante incorpore em seu projeto de vida a participação e a intervenção social no exercício da cidadania, como previsto na \textbf{Competência} Global 6.
Para o desenvolvimento da habilidade EM13LGG201, o componente pode se integrar pelas práticas com usos \textbf{reflexivos} e contextualizados das diversas linguagens ( artísticas, corporais e verbais ), \textbf{mobilizados} por processos comuns à área : \textbf{problematização}, \textbf{análise} e discussão conjunta de como esses usos são marcados pela especificidade dos contextos sociais, culturais e históricos em que se situam.
Dessa forma, a área promove a ampliação dos contextos em que os jovens usam as diferentes linguagens.
O desenvolvimento dessa habilidade evidencia -se na ampliação de repertórios de experiências dos jovens com as linguagens e pela significação delas em seus processos identitários e projetos de vida.
Na habilidade EM13LGG202, o componente pode se integrar pelo procedimento comum de \textbf{análise} de discursos e atos de linguagem — com \textbf{foco} em interesses, relações de poder e perspectivas de mundo — em suas práticas com as linguagens.
O desenvolvimento dessa habilidade evidencia -se em processos de leitura e apreciação \textbf{críticos} e \textbf{reflexivos} sobre esses aspectos, em práticas de linguagem dos diferentes campos de atuação.
A habilidade EM13LGG203 integra -se na qualificação da \textbf{análise} de valores, ideologias e disputas de sentido, nas práticas de linguagem, com \textbf{foco} na \textbf{problematização} de fatores que interferem na valorização e legitimação das práticas, como : por que algumas expressões artísticas são mais valorizadas que as outras?
A quem interessa padronizar uma única variedade linguística como \textbf{legítima}?
Por que determinadas modalidades esportivas recebem mais visibilidade no campo jornalístico-midiático do que outras?
O desenvolvimento dessa habilidade favorece o engajamento dos jovens em favor da legitimação de práticas da cultura juvenil, com uso do que aprenderam sobre \textbf{atores}, processos e dinâmicas da legitimação das práticas.
A habilidade evidencia -se em \textbf{posicionamentos} interessados na diversidade de práticas sociais e engajados na legitimação das práticas de grupos culturais minoritários e/ou tradicionalmente excluídos de reconhecimento social.
O desenvolvimento da habilidade EM13LGG203 favorece a integração ao incorporar em suas práticas \textbf{abordagens} que favoreçam aos estudantes a \textbf{compreensão} da dimensão política das linguagens, com procedimentos de \textbf{análise} de : práticas \textbf{discursivas}, lugares socioculturais que seus sujeitos ocupam e valores e ideias a que se vinculam.
Assim, por meio da \textbf{análise} sistemática dos sentidos e das visões de mundo — que discursos e atos de linguagem expressam, representam e difundem —, a área amplia possibilidades para que os jovens signifiquem essas práticas no (re)conhecimento de si e do outro, valorizando a diversidade em seus processos identitários.
A LI também pode integrar -se à área por propostas de produção de discursos e atos de linguagem que respeitem a diversidade e rompam com padrões de preconceitos, \textbf{materializados} historicamente nas práticas das diferentes linguagens, com protagonismo juvenil na defesa de valores democráticos e dos direitos humanos.
Na oralidade, o desenvolvimento da habilidade EM13LGG301 com integração dentro da própria área é favorecido por propostas de \textbf{autoria} coletiva e colaborativa, com desenvolvimento articulado de dimensões \textbf{relevantes} para a formação integral dos jovens, como : escuta interessada, empatia, diálogo em favor de deliberações e \textbf{consensos}, nas diferentes propostas de experimentação, criação e produção com as linguagens artísticas, corporais e verbais, próprias dos componentes. \textbf{Implica} procedimentos comuns à área : definição e \textbf{análise} de contextos de produção, circulação e \textbf{recepção} de textos e atos de linguagem ; processos de criação, experimentação e produção de textos ; \textbf{problematização} de escolhas de recursos das linguagens e de suas possibilidades de sentidos.
Esta habilidade pode ser desenvolvida no âmbito de projetos da área que oportunizem a \textbf{autoria} coletiva.
Ela favorece o desenvolvimento das \textbf{competências} gerais 4 e 7, por oportunizar o uso \textbf{crítico} das linguagens e o exercício da argumentação, ações estruturantes para que o jovem se posicione no mundo e construa seu projeto de vida, integrando as dimensões pessoal, social/cidadã e profissional.
O desenvolvimento da habilidade EM13LGG302 na integração dentro da própria área pode ser favorecido por procedimentos comuns : \textbf{análise} de textos e atos de linguagem, considerando seus contextos de produção, circulação e \textbf{recepção}, para apreender visões de mundo a que \textbf{remetem} ; e apreciação de ordem ética e estética, com produção de \textbf{posicionamentos} \textbf{críticos}, por meio das diferentes linguagens.
Esse movimento pode ocorrer \textbf{tanto} nas práticas próprias do componente, como em projetos que os integrem, com apreciação de discursos, nos diferentes campos de atuação.
O desenvolvimento dessa habilidade favorece a \textbf{Competência} Global 4, na medida em que o estudante amplia a sua \textbf{capacidade} de utilizar diferentes linguagens para expressar pensamentos, pontos de \textbf{vista} e sentimentos de forma analítica, o que possibilita a construção intencional e informada do projeto de vida.
A habilidade EM13LGG303 pode ser desenvolvida na integração dentro da própria área pelo procedimento comum de \textbf{análise} de diferentes pontos de \textbf{vista}, seus interesses e \textbf{argumentos}, acerca de questões polêmicas de relevância social e de contornos próprios ao repertório de práticas dos diferentes componentes.
Pode também ser favorecida em projetos de investigação e \textbf{debate} de temas controversos de interesse juvenil, com \textbf{mobilização} de performances em que os estudantes se \textbf{invistam} de diferentes \textbf{atores} sociais e exerçam a argumentação, considerando lugares sociais para exercer diferentes pontos de \textbf{vista}.
O desenvolvimento dessa habilidade evidencia -se na manifestação de opiniões sustentadas em \textbf{argumentos}, com respeito a outras perspectivas, ou seja, favorece a \textbf{Competência} Global 7 pelo exercício de \textbf{argumentar} de modo sustentado em fatos, dados e informações que respeitam os direitos humanos e a sustentabilidade \textbf{socioambiental}.
Contribui também para que o estudante faça escolhas informadas e se posicione de modo sustentado, ao longo da construção e concretização de seu projeto de vida.
O desenvolvimento da habilidade EM13LGG304 pode ser favorecido em práticas integradas dos componentes no campo da vida pública, que podem ser organizadas em projetos de intervenção, incluindo os de \textbf{demandas} juvenis, com \textbf{mobilização} do protagonismo dos jovens no \textbf{diagnóstico} de questões que os afetem e/ou afetem a coletividade local e global ; com discussão de possibilidades de intervir por meio das linguagens ; e com usos \textbf{reflexivos} de conhecimentos de gêneros e de práticas da cultura corporal e da Arte.
Seu desenvolvimento evidencia -se no engajamento \textbf{qualificado} de jovens em ações que se pautem pela busca do bem comum, em respeito aos valores democráticos e aos direitos humanos, em consonância com o desenvolvimento da \textbf{Competência} Geral 6.
A área pode favorecer o desenvolvimento da habilidade EM13LGG305, repertoriando estudantes com conhecimento, experimentação e usos dos gêneros e atos de linguagem, para atuar nas diferentes esferas da vida em sociedade, especialmente a política e a \textbf{artístico-cultural}, de modo que possam mapear essas possibilidades de atuação e se valer disso para agir, no exercício de seus \textbf{letramentos}, bem como inovar, valendo -se dos conhecimentos dessas formas para propor outras, com agência inventiva dos jovens no enfrentamento de desafios \textbf{contemporâneos}.
Em conjunto com as outras áreas, é possível desenvolvê -la em propostas de investigação e pesquisa sobre desafios \textbf{contemporâneos}, com complexidade que \textbf{exija} o aporte de diferentes conhecimentos, dentro de projetos que culminem na atuação dos estudantes por meio das linguagens.
Na medida em que \textbf{implica} processos com pensamento \textbf{crítico}, engajamento em favor de mudanças e soluções, o trabalho com essa habilidade pode ser articulado também com os projetos de vida dos estudantes, com \textbf{problematização} de como se veem nessas esferas, o que planejam para si e para os outros, permitindo o desenvolvimento das \textbf{competências} gerais 6 e 10.
Na integração dentro da própria área, o desenvolvimento da habilidade EM13LGG401 pode ser promovido por práticas comuns de \textbf{análise} das condições de produção, circulação e \textbf{recepção} de discursos e atos de linguagem, presentes nas práticas dos diferentes componentes, com \textbf{foco} no reconhecimento da variação linguística, de modo que a área contribua para a generalização da \textbf{compreensão} da língua ( objeto de uso e reflexão comum à área ), como dinâmica, variável, heterogênea e sensível aos diferentes contextos de uso.
A prática comum da \textbf{contextualização} concorre para que o estudante generalize a \textbf{compreensão} das relações entre situação de uso e concretização da língua e que alcance uma visão científica do funcionamento da língua dentro do desenvolvimento da \textbf{Competência} Geral 1, com \textbf{mobilização} desses conhecimentos no combate a preconceitos diversos, sobrepostos ao preconceito linguístico.
O desenvolvimento da habilidade EM13LGG402, na integração dentro da própria área, pode ser favorecido pelos procedimentos comuns de \textbf{contextualização} de diferentes textos e atos de linguagem, em articulação com o desenvolvimento da habilidade EM13LGG401, para a produção de textos e performances, com uso intencional e situado de variedades linguísticas e de estilos das línguas, em combinação com outras linguagens.
O \textbf{intuito} é que, nas práticas próprias dos componentes, haja o exercício comum de um uso plural e \textbf{reflexivo} da língua pelo estudante, com superação de estereótipos e preconceitos que circulam, como o preconceito linguístico.
A etapa de avaliação e revisão dos textos pode ser conduzida dentro dos princípios das metodologias ativas, como oportunidade para que o estudante seja leitor \textbf{crítico} das produções de outros, com \textbf{foco} na \textbf{problematização} de escolhas de variedades e de estilos e \textbf{indicação} de outras possibilidades, dentro dos critérios de adequação e pertinência ao contexto de produção, circulação e \textbf{recepção} do texto.
A habilidade EM13LGG403 integra -se dentro da própria área, pelo uso comum do \textbf{inglês} como língua franca para acessar conhecimentos das práticas de linguagens características dos diferentes componentes, como por exemplo : acessar um tour \textbf{virtual} em um museu, pesquisar no site dos Comitês Internacionais de Esportes, ler jornais em língua \textbf{inglesa} e comparar diferentes visões e interesses que se manifestam na produção de notícias.
Por meio dessas práticas, o estudante desenvolve relações significativas com o \textbf{inglês} em favor da ampliação de seus \textbf{letramentos} e \textbf{multiletramentos}.
Ao vivenciar práticas variadas por meio do \textbf{inglês} e utilizando -se diferentes recursos, inclusive \textbf{digitais}, o jovem pode também considerar essa habilidade de uso no seu \textbf{horizonte} de perspectivas, desejos futuros e planos para seu projeto de vida.
Na integração dentro da própria área, a habilidade EM13LGG501 pode ser favorecida pelos procedimentos comuns de uso e \textbf{análise} de recursos cinésicos ( gestos e expressão corporal ), na apreciação, experimentação e produção de textos e atos de linguagem, dos diferentes componentes, como : uso de gestos e expressões corporais em interações em linguagem oral, em práticas de gêneros como apresentação oral, recital de poesia e leitura de manifestos ; e apreciação de imagens corporais e gestuais, em movimento ou estáticas, como fotografia, artes visuais, \textbf{vídeo}, cinema, dança e teatro.
Os componentes Arte e Educação Física também podem favorecê -la em atividades integradas, dentro de projetos artístico-corporais, especialmente de \textbf{demandas} juvenis, com dança, eixo que é comum a esses componentes.
O desenvolvimento dessa habilidade evidencia -se na exploração \textbf{reflexiva} e intencional de gestos, em diferentes práticas de linguagem, com adequação a contextos e em postura de colaboração com os demais \textbf{atores}.
O desenvolvimento da habilidade EM13LGG502, na integração dentro da própria área, pode ser promovido pela concepção comum de compreender discursos e atos de linguagem como manifestações de ideologias, interesses, preconceitos etc., em articulação com o desenvolvimento de EM13LGG102, operacionalizada na \textbf{análise} de discursos e atos de linguagem que incorram na alimentação de preconceitos e estereótipos em práticas da cultura corporal.
Também pode ocorrer em projetos com \textbf{foco} na investigação de questões \textbf{contemporâneas} que perpassam a discussão de práticas da cultura corporal, em chaves de \textbf{problematização} como : Que estereótipos circulam em relação aos diferentes gêneros no esporte?
A quem interessa mantê -los?
Como combater o racismo na manifestação de torcidas?
Ou, ainda, como apoiar as iniciativas de jovens, com investigação e uso \textbf{crítico} de recursos midiáticos e de meios legais para relatar e se posicionar contra situações de abuso, injustiça e desrespeito nas práticas corporais, em articulação com a habilidade EM13LGG105.
O desenvolvimento dessa habilidade promove o uso consciente e \textbf{crítico} da linguagem corporal, com inserção de valores democráticos na apreciação e uso de práticas corporais no projeto de vida do estudante.
Na integração dentro da própria área e entre as áreas, o desenvolvimento da habilidade EM13LGG503 pode ser promovido pelo envolvimento dos professores e gestores da escola na \textbf{mobilização} dos jovens para iniciativas, planejamento e experimentação de práticas corporais em tempos livres, nos espaços da escola e da comunidade, com protagonismo e autonomia, de modo que possam vivenciar e promover experiências corporais que lhes sejam significativas, ampliando suas possibilidades de vir a ser.
Ela pode ser desenvolvida também em atividades integradas com a área de \textbf{Ciências} da Natureza, na \textbf{problematização} de relações entre recursos usados nas práticas corporais e impactos para o meio e para a saúde, em articulação com o desenvolvimento da habilidade EM13CNT207.
O desenvolvimento dessa habilidade favorece a \textbf{Competência} Global 8, com a estruturação de situações para que o estudante possa se conhecer melhor, cuidar da saúde, integrando aspectos físicos e emocionais ; reconhecer e aprender a lidar com as emoções dos outros e fazer escolhas de práticas corporais no projeto de vida.
Na habilidade EM13LGG601, a integração dentro da própria área pode ser desenvolvida em projetos de investigação e/ou de intervenção acerca do patrimônio artístico local, com engajamento de jovens no conhecimento, na valorização e na difusão dos saberes e fazeres artísticos de sua \textbf{localidade}, em articulação com o desenvolvimento da habilidade de área EM13LGG203.
A habilidade EM13LGG602 integra -se na própria área a partir da consideração de manifestações artísticas diversas nas práticas dos diferentes componentes, com processos de apreciação/fruição pautados por procedimentos comuns de experimentação, \textbf{análise} e \textbf{contextualização}, de modo que temas e reflexões próprias dos diferentes componentes possam ser acionados pelo viés do conhecimento sensível, mesmo que, em conjunto, essas práticas concorram para a ampliação e apropriação de repertórios da Arte.
O acesso a diferentes manifestações artísticas pode se dar pelo uso \textbf{crítico} e ético de novas \textbf{tecnologias} ( exploração de espaços e acervos \textbf{digitais}, participação em canais de coletivos de Arte etc. ), aspectos que se relacionam também ao desenvolvimento da \textbf{Competência} Global 5.
A habilidade EM13LGG603 pode, ainda, ser favorecida pelo desenvolvimento de projetos em que os estudantes tenham protagonismo na iniciativa e experimentação de processos criativos com as diferentes linguagens, sendo possível o desenvolvimento articulado da habilidade de área EM13LGG105, de modo que esses processos sejam também tomados como formas de intervir em realidades locais ou globais — nesse caso, em combinação com formas de difusão das produções mediadas pelas novas \textbf{tecnologias}, em articulação com o desenvolvimento da habilidade EM13LP17.
Ela pode também ser desenvolvida no âmbito das práticas \textbf{sugeridas} para o desenvolvimento da habilidade EM13LP47 ( mostras, saraus, batalhas, oficinas ), tomando -se a participação nesses \textbf{eventos} como contexto dos processos de criação com as diferentes linguagens artísticas.
Ao prever a convergência de conhecimentos e experiências na utilização de diversas linguagens artísticas, a habilidade promove o desenvolvimento da \textbf{Competência} Global 4.
Na habilidade EM13LGG604 a integração dentro da própria área pode ser favorecida em articulação com o desenvolvimento da habilidade de área EM13LGG602 e no bojo das práticas \textbf{sugeridas} para ela, \textbf{implicando} processos de \textbf{análise} dos contextos de produção, circulação e \textbf{recepção} das produções escolhidas, com \textbf{foco} no exercício de \textbf{contextualização} nas práticas artísticas de uma época e com relações a aspectos sociais, políticos, \textbf{econômicos} e culturais.
A habilidade EM13LGG701 pode ser desenvolvida pelo exercício comum de considerar as práticas mediadas pelas \textbf{tecnologias} \textbf{digitais} da informação e comunicação, nas práticas de linguagem próprias de cada componente, com procedimentos de experimentação, \textbf{análise} e \textbf{problematização} de princípios e valores nos usos.
É possível também promover o desenvolvimento dessa habilidade por meio de projetos da área e entre as áreas que permitam ao estudante trabalhar princípios, funcionalidades e uso ético, criativo e responsável das TDIC, em práticas autorais e coletivas e em diálogo com práticas das culturas juvenis.
O desenvolvimento da habilidade EM13LGG702 pode ser promovido pelos procedimentos comuns de qualificação dos usos e mecanismos de investigação, pesquisa e curadoria de informação e opinião, nas práticas dos diferentes componentes.
Seu desenvolvimento também pode ser favorecido no âmbito de projetos, especialmente por \textbf{demandas} juvenis, que \textbf{impliquem} práticas de linguagem em ambiente \textbf{digital} e concorram para os processos formativos da comunidade escolar e/ou de outros sujeitos locais e globais, como exercício ético e \textbf{crítico} na produção de textos e atos de linguagem, considerando -se os meios de produção, circulação e \textbf{recepção}.
O desenvolvimento da habilidade EM13LGG703, integrada à própria área, pode ser promovida pelos componentes em projetos que deem centralidade aos jovens na \textbf{autoria} coletiva, em processos de criação e produção de textos e atos de linguagens que \textbf{impliquem} : uso \textbf{reflexivo} de ferramentas, softwares e mídias \textbf{digitais} ; \textbf{análise}, experimentação, combinação e \textbf{edição} de recursos semióticos variados ; articulação de mídias ( multimídia, transmídia etc. ) ; criatividade ; colaboração ; diálogos com práticas das culturas juvenis e projetos de vida ; e reflexão sobre a necessidade da promoção de valores éticos e não proliferação de discursos de ódio.
Observa -se a possibilidade de desenvolver essa habilidade em articulação com EM13LGG105 e EM13LGG304, no bojo das práticas \textbf{sugeridas}.
Por fim, a habilidade EM13LGG704 integra -se dentro da própria área por meio das práticas de investigação e pesquisa dos diferentes componentes, com procedimentos comuns de uso \textbf{crítico} e \textbf{reflexivo} de recursos \textbf{digitais} de seleção, filtragem, checagem, validação, comparações, \textbf{análises}, (re)organização, categorização, reedição de informações ou, ainda, no âmbito de projetos da área que \textbf{impliquem} curadoria e/ou redistribuição de informação/opinião ( Reúna, 2020 ).
Dentre as considerações \textbf{teóricas} e os procedimentos que possam \textbf{subsidiar} as práticas de linguagem, abrem -se para contribuições as várias possibilidades de \textbf{teorias} linguísticas capazes de suscitar reflexões, conscientização linguística e \textbf{compreensão} de como a língua é usada pelos sujeitos – em seus discursos, em situações sociais de uso ( na vida real ) e por meio da \textbf{compreensão} das nuances do explícito e implícito, inferido e/ou compreendido ideologicamente.
Isso pode ser obtido por meio da Análise do Discurso, pela Linguística Textual e seus processos, pela Pragmática, para a \textbf{compreensão} dos significados em contextos de interações e atos de fala, pela Fonética, em favorecimento da \textbf{compreensão} de sons e pronúncias na diversidade de usuários da língua ; pela Semântica, com significados das palavras e relações aos ditos em textos ; pela Sintaxe, na formação de enunciados ; pela Morfologia, com a formação das palavras e seus significados no repertório linguístico da língua utilizada ; pela Semiótica, seus recursos e a \textbf{compreensão} de sinais, signos e os sentidos por eles produzidos ; pela Sociolinguística, para a \textbf{compreensão} das relações entre o uso da linguagem e a sociedade, por meio de possibilidades que podem \textbf{influenciar} ou determinar escolhas linguísticas e/ou o comportamento e uso da língua ; pelo Interacionismo Sociodiscursivo ( ISD ) e os \textbf{subsídios} \textbf{teóricos} que \textbf{embasam} o trabalho com gêneros textuais e \textbf{discursivos}, favorecendo aos estudantes o desenvolvimento de habilidades de \textbf{compreensão}, \textbf{recepção} e produção textual, com possibilidades de maior \textbf{aprofundamento}, pelos processos de \textbf{contextualização} utilizados na \textbf{abordagem} das práticas \textbf{discursivas} ( Bronckart, 1999, 2003, 2007, 2009 ; Cristóvão, 2007, 2008 ).
Em consonância com as perspectivas de desenvolvimento de \textbf{competências} e habilidades no ensino de LI, os pressupostos \textbf{teórico-metodológicos} que \textbf{subsidiam} esse documento curricular orientam para a \textbf{abordagem} das atividades na realização de “ leitura, escuta, apreciação, experimentação e \textbf{análise} de discursos e atos de linguagem ” com base na \textbf{contextualização}, \textbf{problematização} e investigação ( BNCC/Reúna, 2020 ), \textbf{aproximando}, assim, o estudante do contexto significativo para sua \textbf{recepção}, \textbf{compreensão} e produção ( oral e escrita ), em suas práticas de linguagem.
Para \textbf{tanto}, em benefício da organização de procedimentos que possam contribuir : a ) para o desenvolvimento de leituras \textbf{críticas} na \textbf{compreensão} leitora ; b ) para o desenvolvimento da oralidade, produção textual e expressão oral ( réplicas, performances ) ; c ) para a produção textual e expressão escrita ; d ) para a conscientização e \textbf{análises} linguísticas e semióticas ; e ) para a conscientização \textbf{intercultural} e a interculturalidade, acrescentamos considerações \textbf{teóricas} e sugestões de \textbf{encaminhamentos} pedagógicos, como as apresentadas a seguir : a ) Para práticas \textbf{discursivas} de leitura e \textbf{compreensão} leitora, considerando o estudante como centro da aprendizagem, na perspectiva de seu engajamento \textbf{crítico}, no ato de ler, e para que sua resposta ao texto seja alcançada, será necessário \textbf{mobilizar} seus conhecimentos e conduzi -lo às inferências a respeito dos significados e sentidos encontrados por ele, pela \textbf{contextualização}.
Assim, \textbf{sugere} -se considerar etapas de \textbf{compreensão} leitora envolvidas em responder criticamente ao texto, como : a leitura textual inicial e a \textbf{compreensão} do contexto de produção ; fazer a leitura interna do texto buscando a \textbf{compreensão} da lógica dos \textbf{argumentos} e evidências apresentadas pelo \textbf{autor}, fatos, intenções, ideologias e \textbf{maneiras} de \textbf{dizer} o que é \textbf{dito} ; ler para além do texto, por meio de questionamentos, e comparando o que se sabe a respeito do texto e o que já foi \textbf{dito} por outros ; realizar o contraponto da leitura avançando com questões de \textbf{problematização} e contrastando com o que outros já \textbf{disseram} — o que requer múltiplas leituras, a investigação, para que possa seguir questionando, buscando a \textbf{compreensão}, e o desenvolvimento de estratégias de leitura e de repertório linguístico, significando e fazendo sentido em textos, enquanto busca por respostas ( Beach ; Webb, 2016 ) e constrói conhecimentos e \textbf{competências} leitoras ( Soares, 2020 ).
Contextualizações para a \textbf{compreensão} leitora podem ser otimizadas pela integração da leitura e oralidade, favorecendo os estudantes em suas interações com o texto por meio da dialogia oportunizada por questionamentos e respostas sobre o contexto e a produção textual, ou seja, à medida que se lê, se analisa e se comentam os significados.
Considerados esses \textbf{parâmetros}, pode -se ampliar a \textbf{compreensão} pelo estabelecimento de hipóteses sobre o conteúdo temático, sobre o gênero \textbf{discursivo}, o qual \textbf{materializa} o texto num determinado suporte e meio de circulação social, sua finalidade, objetivos e intencionalidades, e reconstruir conhecimentos pela reflexão, \textbf{crítica} e investigação.
Favorecendo atividades para possibilitar os \textbf{multiletramentos}, apontamos, à leitura de textos \textbf{digitais}, as vastas possibilidades de \textbf{contextualizações}, verbais e não verbais, dadas ao aporte possível de diferentes linguagens que compõem o texto além das palavras, como a não linearidade textual – colunas, imagens coloridas, gráficos, sons, diagramas, símbolos, signos, links, entre outros -, e a linguagem “ híbrida, dinâmica e flexível, que dialoga com outras \textbf{interfaces} semióticas, adiciona e acondiciona a sua superfície formas outras de textualidade ” ( Xavier, 2005, p. 171 ), de modo que a interatividade do estudante-leitor se torna ampliada.
Nas práticas \textbf{discursivas} de leitura, o engajamento \textbf{crítico} é favorecido por meio das réplicas possibilitadas pelas discussões com base na reflexão \textbf{crítica} sobre o texto lido. “ \textbf{Posicionamentos} responsáveis em relação a temas, visões de mundo e ideologias veiculados por textos e atos de linguagem ” ( Reúna ; BNCC, 2020 ) são encorajados, não somente pelo texto em si, mas incluindo considerações dos contextos histórico, cultural, social ou político, no qual o texto foi produzido e deve ser compreendido pelos estudantes ( Beach ; Webb, 2016, p.104 ).
A \textbf{compreensão} e a \textbf{capacidade} de interpretação dos estudantes podem ser ampliadas ao se permitir que discutam livremente suas apreciações, impressões e pontos de \textbf{vista} ( Inglês-Português pode ser necessário ).
Marca -se nessa \textbf{instância} a importância da transversalidade da dimensão cultural e \textbf{intercultural} encontrada nas temáticas abordadas a partir dos textos provindos de diferentes campos de atuação e esferas de circulação social, de diferentes sociedades e culturas.
Nessa mesma dimensão, a \textbf{compreensão} leitora de textos literários é responsável em estimular a imaginação dos estudantes, levantando questionamentos sobre relações humanas, sociais, históricas e culturais.
Uma das contribuições mais significativas da leitura literária \textbf{crítica} está no desenvolvimento de empatia e tolerância moral.
Em resposta a textos literários está a oportunidade do desenvolvimento da empatia com as experiências dos personagens e \textbf{autores}, seus pensamentos e emoções, passíveis de trazer \textbf{influências} de atitudes e crenças, o que \textbf{influencia}, por sua vez, as percepções de \textbf{eventos} ou representatividades sociais, levanta reflexões e réplicas substancialmente \textbf{críticas} e engajadas com as múltiplas \textbf{vozes}, e traz oportunidades para o exercício de mudanças de atitudes de preconceito em relação ao outro em desvantagem ( Beach ; Webb, 2016, p. 99 ).
Às práticas \textbf{discursivas} de linguagem, esse referencial curricular de LI acrescenta, com destaque, o estabelecido pelas Leis nas Diretrizes e Bases da Educação Nacional, para a inclusão no currículo oficial da rede de ensino a obrigatoriedade das temáticas : História e Cultura Afro-Brasileira ; História e Cultura Afro-brasileira e Indígena ; Relações Étnico-Raciais e para o Ensino de História e Cultura Afro-Brasileira e Africana ( Brasil, 2018, p. 20 ). b ) As práticas de \textbf{compreensão} e produção oral, para o desenvolvimento das habilidades e \textbf{competências} específicas das linguagens na oralidade ( ILF ), são integrantes e \textbf{imprescindíveis} no processo de comunicação.
A fala, a escuta e audição são \textbf{fundamentos} do \textbf{letramento} em todas as suas dimensões.
Estratégias e práticas \textbf{discursivas} de oralidade contribuem para o desenvolvimento da \textbf{capacidade} de expressão oral, da socialização e da autonomia, bem como para o pensamento, a sensibilidade ética, estética, política e cultural, e para o enriquecimento da construção de conhecimentos.

% matched lemmas: abordagem, aluno, análise, aprimorar, aprofundamento, aproximar, argumentar, argumento, artístico-cultural, ator, aula, autor, autoria, capacidade, ciência, coesão, competência, complexo, compreensão, consenso, contemporâneo, contextualização, crítico, debate, demanda, descrição, diagnóstico, digital, discursivo, dizer, econômico, edição, embasar, encaminhamento, ensino-aprendizagem, entrevista, evento, exigir, foco, fundamento, grifo, horizonte, ideal, impactar, implicar, imprescindível, indicação, influenciar, influência, inglês, instância, intercultural, interdisciplinaridade, interface, intuito, investir, legítimo, letramento, localidade, maneira, materializar, metodológico, mobilizar, mobilização, multiletramento, multissemiótico, organizador, parâmetro, permear, pluralidade, posicionamento, problematização, proficiência, progressão, quadro, qualificar, recepção, reflexivo, reiterar, relevante, remeter, sequência, sistematizar, socioambiental, subsidiar, subsídio, sugerir, tanto, tecnologia, teoria, teórico, teórico-metodológico, virtual, vista, voz, vídeo
\end{textsample}
